\tocsection{\textbf{11/16/26 Lecture 33:} $H_1$ and the Fundamental Group, Simplicial Maps}

\begin{definition}[\textbf{Definition 33.1:} Elementary 1-Cycles]
    Let $K$ be a simplicial complex. An \emph{elementary $1$-cycle} is the cycle coming from a simple (non self-intersecting) polygonal loop in $K$.
\end{definition}

\begin{lemma}[\textbf{Fact 33.2}]
    $Z_1(K)$ is generated by elementary $1$-cycles.
\end{lemma}

\begin{theorem}[\textbf{Theorem 33.3}: $H_1(K) = \pi_1^{ab}(\abs{K}, v)$]
    If $K$ is connected, $H_1(K)$ is the abelianization of $\pi_1(\abs{K}, v)$.
\end{theorem}
\begin{proof}{Proof:}
    Given an edge loop $\alpha = vv_1 \cdots v_kv$, we get a $1$-cycle, omitting the edge $(v_iv_{i + 1})$ if $v_i = v_{i + 1}$, by
    \[ z(\alpha) = (vv_1) + \cdots + (v_kv) \in Z_1(K) \]
    Thus, $z$ sends edge loops at $v$ to $Z_1(K)$.\\

    Now, $E(K, v)$ consists of edge loops based at $v$ modulo the moves
    \fig{0.6}{Lecture33Fig1}{}
    and $(vv) = e$. We saw $E(K, v) \cong \pi_1(\abs{K}, v)$.\\

   
    Now, we get a homomorphism $\phi : E(K, v) \to H_1(K)$ by sending $\alpha$ to
    \[ \phi(\{\alpha\}) = z(\alpha) \bmod B_1(K) \]

    It can be checked that $\phi$ is well defined. We want to show $\phi$ is onto and $\ker(\phi) = [E(K, v), E(K, v)]$, which will show
    \[ H_1(K) \cong E(K, v)/[E(K, v), E(K, v)] \cong \pi_1^{ab}(\abs{K}, v) \]


    We need to show $\phi$ is onto. Suppose $z_1$ is an elementary $1$-cycle,
    \[ z_1 = (w_1w_2) + \cdots + (w_sw_1) \]
    Choosing an edge path $\gamma$ from $v$ to $w_1$, let
    \[ \alpha = \gamma w_1 \cdots w_sw_1\gamma^{-1} \]
    Thus, $z(\alpha) = z_1$ via cancellation. Then, as the elementary $1$-cycles generate $Z_1(K)$, their cosets (mod $B_1(K)$) generate $H_1(K)$.\\

    Now, we need to show $\ker(\phi) = [E(K, v), E(K, v)]$. Now, as
    \begin{align*}
        \phi(g_ig_i'g_i^{-1}(g_i')^{-1}) &= \phi(g_i)\phi(g_i')\phi(g_i)^{-1}\phi(g_i')^{-1}\\
        &= e
    \end{align*}
    since $H_1(K)$ is abelian, we see that $[E(K, v), E(K, v)] \subseteq \ker(\phi)$.\\
    
    Now, we want to show $\ker(\phi) \subseteq [E(K, v), E(K, v)]$. Suppose $\alpha$ is an edge loop such that $\{\alpha\} \in \ker(\phi)$. Thus, $\phi(\{\alpha\}) = 0$ in $H_1(K)$, so $z(\alpha) \in B_1(K)$.\\

    Writing $\alpha = vv_1 \cdots v_kv$, write
    \[ z(\alpha) = \partial(\lambda_1\sigma_1 + \cdots + \lambda_{\ell}\sigma_{\ell}) \]
    where $\{\sigma_i\}$ are oriented $2$-simplices in $K$.\\

    Write $\sigma_i = (a_ib_ic_i)$. Now, we see
    \fig{0.75}{Lecture33Fig2}{}
    where $\gamma_i$ is an edge path from $v$ to $a_i$. Looking at the edge loop $\gamma_ia_ib_ic_ia_i\gamma_i^{-1}$, by our edge moves, $\gamma_ia_ib_ic_ia_i\gamma_i^{-1}$ is trivial in $E(K, v)$.\\

    Let
    \[ \beta = \prod_i (\gamma_ia_ib_ic_ia_i\gamma_i^{-1})^{\lambda_i} \]
    Since each of the factors is trivial, $\beta$ is trivial.\\

    Thus, $\{\alpha\} = \{\alpha\beta^{-1}\}$ in $E(K, v)$. Now,
    \begin{align*}
        z(\gamma_ia_ib_ic_ia_i\gamma_i^{-1}) &= z(a_ib_ic_ia_i) \quad \tt{(paths cancel going forward and backward)}\\
        &= (a_ib_i) + (b_ic_i) + (c_ia_i)\\
        &= \partial(\sigma_i)
    \end{align*}
    Thus, $z(\alpha\beta^{-1}) = 0$ as a $1$-cycle.
    Now, $\alpha\beta^{-1}$ is an edge loop based at $v$. Thus, every oriented edge $(ab)$ appears in $\alpha\beta^{-1}$ as many times as $(ba)$.\\
.
    Now, recall that if $L$ is a maximal tree, we get a group $G(K, L)$ and an isomorphism $\theta : E(K, v) \to G(K, L)$. Recall $G(K, L)$ is generated by the edges not in $L$. Thus, by above, we see
    \[ \theta(\{\alpha\beta^{-1}\}) \in [G(K, L), G(K, L)] \]
    so as $\theta$ is an isomorphism, we see $\{\alpha\beta^{-1}\} \in [E(K, v), E(K, v)]$, so $\{\alpha\} \in [E(K, v), E(K, v)]$. Thus, we have everything we wanted to show.

    $\qed$
\end{proof}

\begin{example}{\textbf{Example:} The Torus}
    \fig{0.6}{Lecture33Fig3}{}
    Since $\pi_1(T^2) = \langle a, b \mid aba^{-1}b^{-1} \rangle = \Z^2$ is abelian, $H_1(T^2) = \Z^2$.
\end{example}

\begin{example}{\textbf{Example:} The Double Torus}
    \fig{0.55}{Lecture33Fig4}{}
    We saw
    \[ \pi_1(T_2^2) = \langle a_1, b_1, a_2, b_2 \mid a_1b_1a_1^{-1}b_1^{-1}a_2b_2a_2^{-1}b_2^{-1} \rangle \]
    So abelianizing, we see $H_1(T_2^2) = \Z^4$.
\end{example}

In general, if $\Sigma_g$ is the closed orientable surface of genus $g$:
\fig{0.6}{Lecture33Fig5}{}
\begin{center}
    (left to right: $g = 0$, $g = 1$, $g = 2$)
\end{center}
We have
\[ H_1(\Sigma_g) \cong \Z^{2g} \]

\subsection*{Maps}
\sectionspace

Let $K, L$ be simplicial complexes and let $s : \abs{K} \to \abs{L}$ be a \emph{simplicial map}; that is, $s$ is continuous and sends simplices linearly to simplices.\\

Then, we want to define $s_{q*} : H_q(K) \to H_q(L)$.\\

First, define $s_q : C_q(K) \to C_q(L)$ by sending an oriented $q$-simplex $\sigma = (v_0 \cdots v_q)$ to
\[ s_q(\sigma) = \left\{\begin{matrix} (s(v_0) \cdots s(v_q)) & \tt{if the } s(v_i) \tt{ are all distinct} \\ 0 & \tt{else} \end{matrix}\right. \]

We will prove the following lemma next time.

\begin{lemma}[\textbf{Lemma 33.4}]
   
    $\partial \circ s_q = s_{q - 1} \circ \partial$ as maps $C_q(K) \to C_{q - 1}(L)$. That is, the following diagram commutes. 
    \begin{center}
        \begin{tikzcd}
            C_q(K) \arrow[r, "\partial"] \arrow[d, "s_q"'] & C_{q - 1}(K) \arrow[d, "s_{q - 1}"] \\
            C_q(L) \arrow[r, "\partial"] & C_{q - 1}(L)
        \end{tikzcd}
    \end{center}
\end{lemma}
